% !TeX spellcheck = en_US
\addsection{Deck-building}{\skills/wisdom.png}

\begin{multicols}{2}

\subsection*{\pagetarget{Playerdecks}{Player Decks}}
All players have a unique Deck which represents their Main Hero's Abilities and Equipment.
Decks may contain Statistic, Ability, Spell, Artifact and the Main Hero's Specialty cards.
Each player's Deck starts with 9 cards, built during the game's setup.\par
\bigskip
\bigskip

\subheader{General Card Rules}
\begin{enumerate}
  \item Cards can be played \textbf{only on your Turn}, or in a \pagelink{Combat}{Combat} involving your \textbf{Main Hero}.
  \item After a card is used, discard it.
    Each player has their own separate Discard Pile\index{Discard Pile}.
  \item If your Deck is empty when you need to draw a card, \textbf{shuffle your Discard Pile} into a new Deck to draw from.
  \item Whenever your Hero gains a card for any reason, put it \textbf{directly into your hand} unless otherwise stated.
  \item Whenever you are instructed to \textbf{Search} (X)\index{Search (X)} the Ability, Artifact, or Spell Deck, you may either look at the top (X) cards from the specified Deck, take one of them to your hand, and discard the others, \textbf{OR} instead of looking at the top (X) cards, gain the top card from that Deck's Discard Pile.
  \item The Ability, Artifact, and Spell Decks each have their own Discard Piles, created during the setup, which help you identify these Decks.
    If a Deck ever runs out of cards, reshuffle it and discard its top card to form a new Discard Pile.
    Whenever one of these Discard Piles is empty, \textbf{refill it} with that Deck's top card.
  \columnbreak
  \item Cards have the following \pagetarget{Card Effects}{types of effects}:
  \begin{itemize}
    \item \textbf{Instant} \svg{instant} Effects are resolved immediately.
    \item \textbf{Activation} \svg{activation} Effects must be played when Activating your own unit in Combat.
    \item \textbf{Map} \svg{map} Effects cannot be used during Combat.
    \item \textbf{Ongoing} \svg{ongoing} Effects\index{Ongoing Effects} last until they are used up or until the player who played them starts their next Turn (whichever happens first).
      \begin{expansion}[before skip balanced=0.3em,]{stronghold,conflux,cove}
      \pagetarget{Time Token}{}
      \begin{tikzpicture}[overlay]
        \node at (5.5, 0) {
\includegraphics[width=0.22\linewidth]{\images/time-token.png}};
      \end{tikzpicture}\parbox{0.75\hsize}{You can use Time Tokens\index{Time Token} to track the duration of Ongoing Effects.}
      \end{expansion}
    \item \textbf{Permanent} \svg[7]{permanent} cards\index{Permanent cards} stay in play until discarded or replaced.
      They are played the same way as \svg{activation} and \svg{map} cards.
      \textbf{You may only have one permanent card at a time}; playing another discards the first.
      While a Permanent card is in play, you may use either its Basic or Expert Effect, but not both simultaneously.
  \end{itemize}
  \item Whenever you are instructed to \textbf{remove a card}, remove it completely from the game.
\end{enumerate}

\vfill
\hfill{\transparent{0.2}
\includegraphics[width=0.9\linewidth]{\art/quicksand.png}}
\end{multicols}
\clearpage

\subheader{Searching Example}
\begin{multicols*}{2}

    \textit{Bob, playing as Deemer the Warlock, has built a Mage Guild during the previous Game Round, enabling him to now use his Spell Book Token to purchase additional Spells.}

    \includegraphics[width=\linewidth]{\examples/deemer_receiving_spellbook.png}\par
    \textit{He spends the token and pays the cost of 5 \svg{gold} indicated on his Town Board, allowing him to \textbf{Search (2)} the Spell Deck.}\par
    \includegraphics[width=\linewidth]{\examples/deemer_spell_cost.png}\par
    \filbreak
    \textit{He now has the choice of either gaining the top Card of the Spell Card Discard Pile, or drawing the top two cards of the Spell Deck and gaining one of them.
        He is not interested in gaining the Spell Curse, which is on top of the Discard Pile, and instead decides to draw two cards – a Magic Arrow and a Fireball.}
    \includegraphics[width=\linewidth]{\examples/fireball_over_magic_arrow.png}\par
    \textit{He decides to keep the Fireball, placing it into his hand and discarding the Magic Arrow into the Spell Discard Pile.}

    \vspace*{\fill}

    \hfill{\transparent{0.2}
\includegraphics[width=0.8\linewidth]{\art/fire_magic.png}}

\end{multicols*}

\pagetarget{Statistic}{}\pagetarget{Ability}{}
\subheader{Statistic and Ability Cards}
\begin{multicols}{2}

All Statistic\index{Statistic cards} and Ability cards\index{Ability card} have a \textbf{Basic Effect}\index{Basic Effect} and a stronger \textbf{Expert} \svg{expert} \textbf{Effect}\index{Expert Effect}, which is shown below the Basic Effect.
Whenever you play a Statistic or Ability card, you must choose which effect you are using.
The number of \svg{expert} Effects you can use each Round is limited by your Main Hero's \pagelink{Level}{Level}.
Track the number of uses you have in any suitable manner, such as by moving Black Cubes on and off your Hero Board.\par

\columnbreak

\note{11}{
  \wikilink{abilities/necromancy}{Necromancy cards} are limited to the Necropolis Faction \svg{necro-yellow}.
  When a non-Necropolis player draws a Necropolis card\index{Necropolis card} from the Ability Deck, they may either discard it and draw a new card as a replacement or gain it.
  Non-Necropolis players \textbf{cannot use} Faction Specific cards from their hand in any way besides for effects that discard them.
}
\end{multicols}

\columnratio{0.25}
\begin{paracol}{2}
  \begin{expansion}[title=\phantom{\textbf{\tiny{foo}}}, overlay={}]{}
    \begin{center}
      \begin{tikzpicture}
        \node (card) {\includegraphics[width=3.5cm]{\cards/knowledge.png}};
        \begin{relcoords}{card}
          \draw (0.65, 0.6) node {\encircle{1}};
          \draw (0.28, 0.4) node {\encircle{2}};
          \draw (0.28, 0.2) node {\encircle{3}};
        \end{relcoords}
      \end{tikzpicture}
      \imagecaption{Statistic card\\\phantom{Blah}}
    \end{center}
  \end{expansion}
  \switchcolumn
  \begin{expansion}{inferno}
    \hspace*{-2.3mm}
    \begin{minipage}[t]{0.3\linewidth}
      \vspace*{0mm}
      \begin{center}
        \begin{tikzpicture}
          \node (card) {\includegraphics[width=3.5cm]{\cards/empowered-knowledge.png}};
          \begin{relcoords}{card}
            \draw (0.65, 0.8) node {\encircle{1}};
            \draw (0.28, 0.25) node {\encircle{2}};
          \end{relcoords}
        \end{tikzpicture}
        \imagecaption{Empowered\\Statistic card}
      \end{center}
    \end{minipage}
    \hfill
    \begin{minipage}[t]{0.65\linewidth}
      \vspace*{0mm}
      \subsection*{\pagetarget{Empowered Statistic}{Empowered} Statistic cards}\index{Empowered Statistic Card}
      These cards are more powerful versions of the normal Statistic cards.
      They have only one effect which is identical to the normal Statistic's Expert effect, but does not require using your \svg{expert}.\par

      \medskip
      These cards can be gained via certain game effects, such as Visiting a \pagelink{Star Axis}{Star Axis}.
    \end{minipage}
  \end{expansion}
\end{paracol}
\vspace*{2em}
\begin{paracol}{2}
  \begin{expansion}[title=\phantom{\textbf{\tiny{foo}}}, overlay={}]{}
    \begin{center}
      \begin{tikzpicture}
        \node (card) {\includegraphics[width=3.5cm]{\cards/necromancy_card.png}};
        \begin{relcoords}{card}
          \draw (0.65, 0.6) node {\encircle{1}};
          \draw (0.28, 0.4) node {\encircle{2}};
          \draw (0.28, 0.2) node {\encircle{3}};
          \draw (0.9, 0.1) node {\encircle{4}};
          \draw (0.9, 0.3) node {\encircle{4}};
        \end{relcoords}
      \end{tikzpicture}
      \imagecaption{Ability card\\\phantom{Blah}}
    \end{center}
  \end{expansion}
  \switchcolumn
  \begin{expansion}{navalbattles}
    \hspace*{-2.3mm}
    \begin{minipage}[t]{0.3\linewidth}
      \vspace*{0mm}
      \begin{center}
        \begin{tikzpicture}
          \node (card) {\includegraphics[width=3.5cm]{\cards/empowered-ability.png}};
          \begin{relcoords}{card}
            \draw (0.65, 0.8) node {\encircle{1}};
            \draw (0.28, 0.25) node {\encircle{2}};
            \draw (0.9, 0.1) node {\encircle{4}};
          \end{relcoords}
        \end{tikzpicture}
        \imagecaption{Empowered\\Ability card}
      \end{center}
    \end{minipage}
    \hfill
    \begin{minipage}[t]{0.65\linewidth}
      \subsection*{\pagetarget{Empowered Ability}{Empowered} Ability cards}\index{Empowered Statistic card}
      The Empowered Ability cards work very similarly to Empowered Statistic cards.
      Their only effect corresponds to the Expert effect of normal Ability cards, but without requiring you to use your \svg{expert}.

      \medskip
      These cards can be gained by spending \pagelink{Empowered Ability Token}{Empowered Ability Tokens}, which are only available in certain Scenarios (e.~g.~``Island of Fire'' from the Naval Battles Mission Book).
    \end{minipage}
  \end{expansion}
\end{paracol}

\bigskip
\begin{multicols}{2}
\begin{multicols*}{2}
  \footnotesize
  \begin{itemize}
    \item[\textbf{1.}] Name
    \item[\textbf{2.}] Basic Effect\\
    \item[\textbf{3.}] Expert Effect
    \item[\textbf{4.}] Necropolis-specific card
  \end{itemize}
\end{multicols*}
\columnbreak
\centering
\end{multicols}

\clearpage
\begin{multicols}{2}
\pagetarget{spells}{\subheader{Spell Cards}}

Spell cards\index{Spell card} have three possible primary effects.
Using the topmost, basic version of the Spell usually has no additional costs.
To access the other effects, you may \textbf{Empower} a Spell\index{Empowering Spells} by paying the indicated cost (3) to get a more powerful outcome (4).
You may pay this cost by playing other cards for their Empower \svg{empower} effect (e.g. the \wikilink{statistics/power}{Power} Statistic card) before casting the Spell.
All Spell cards also have an alternative bottom (5) \svg{empower} effect.
You cannot preserve the played \svg{empower}: casting a Spell consumes all the \svg{empower} played so far, and if you stop playing cards, the power dissipates.
You may use a cheaper \svg{empower} effect than what you've paid for, though.\par

\bigskip
\begin{center}
  \begin{scriptsize}
    \begin{tikzpicture}
      \node (card) {\includegraphics[width=0.6\linewidth]{\cards/spell.png}};
      \begin{relcoords}{card}
        \draw (0.25, 0.85) node {\encircle{1}};
        \draw (0.2, 0.6) node {\encircle{2}};
        \draw (0.32, 0.35) node {\encircle{3}};
        \draw (0.66, 0.35) node {\encircle{4}};
        \draw (0.62, 0.1) node {\encircle{5}};
      \end{relcoords}
    \end{tikzpicture}
  \end{scriptsize}\\
  \imagecaption{Spell card}\\
\end{center}

\begin{multicols*}{2}
  \footnotesize
  \begin{itemize}[itemsep=0pt, parsep=5pt, topsep=0pt, partopsep=0pt]
    \item[\textbf{1.}] Spell Name
    \item[\textbf{2.}] School of Magic
    \item[\textbf{3.}] Cost to\\Empower
    \columnbreak
    \item[\textbf{4.}] Spell Effect
    \item[\textbf{5.}] Alternative\\Effect
  \end{itemize}
\end{multicols*}

Spells can be gained by Building the Mage Guild\index{Mage Guild}.
Building it also enables you to use the Spell Book Token\index{Spell Book Token} for buying additional Spells.
The Token cannot be used during the \textbf{same Round} when the Guild is constructed. Spells can also be \pagelink{Trading}{traded} in Alliance and Cooperative Scenarios.

During Combat, each player may play \textbf{only one Spell card per Combat Round}.

\columnbreak
When you use an effect which lets you cast a Spell from somewhere other than your hand (such as directly from the Spell Deck), place that Spell into your Discard Pile after casting it.

Some complex Spell cards use keywords or require special components, see \pagelink{spell keywords}{Card Effect Keywords}.

\begin{center}
  \vspace*{\fill}
  {\transparent{0.2}
\includegraphics[width=0.9\linewidth]{\art/air_magic.png}}
  \vspace*{\fill}
\end{center}

\subsection*{\pagetarget{Schools of Magic}{Schools of Magic}}
All Spell cards belong to one School of Magic\index{Schools of Magic}: either Air, Fire, Earth or Water.
Spells with one School symbol on them are called \textbf{Basic Spells}\index{Basic Spell}, while Spells with four identical symbols are called \textbf{Expert Spells}\index{Expert Spell}.
\textbf{\wikilink{spells/magic_arrow}{Magic Arrow}} is a Basic Spell; when casting the \textbf{Magic Arrow}, you must select which School it belongs to.
By default there's no difference, but certain game effects can refer to specific Schools and to Basic/Expert Spells.

\begin{minipage}[t]{0.48\textwidth}
  \centering
  \minipage[b]{0.5\textwidth}
    \centering
    
\includegraphics[width=0.75\linewidth]{\images/school_of_fire.png}
    \imagecaption{School of Fire}
  \endminipage
  \minipage[b]{0.5\textwidth}
    \centering
    
\includegraphics[width=0.75\linewidth]{\images/school_of_water.png}
    \imagecaption{School of Water}
  \endminipage
  \hfill\allowbreak%
  \bigbreak
  \minipage[b]{0.5\textwidth}
    \centering
    
\includegraphics[width=0.75\linewidth]{\images/school_of_air.png}
    \imagecaption{School of Air}
  \endminipage
  \minipage[b]{0.5\textwidth}
    \centering
    
\includegraphics[width=0.75\linewidth]{\images/school_of_earth.png}
    \imagecaption{School of Earth}
  \endminipage
  \bigbreak
\end{minipage}
\clearpage

\pagebreak
\pagetarget{Specialty}{\subheader{Hero Specialty Cards}}

Hero Specialty cards\index{Specialty card} are gained from \pagelink{Level}{Level ups}.
Each Main Hero has a unique set of Specialty cards.
While some of these cards have effects which resemble Spell cards and even make use of \svg{empower}, Specialty cards are their \textbf{own unique category of cards}.
For example, the limit of 1 Spell per Combat Round doesn't apply to Specialty cards.

{
    \medskip
    \centering
    \includegraphics[width=0.7\linewidth]{\cards/specialty.png}\\
    \medskip
    \footnotesize\textit{A Level IV Specialty card, belonging to Catherine the Knight.}
}

{
    \medskip
    \centering
    \includegraphics[width=0.7\linewidth]{\cards/specialty2.png}\\
    \medskip
    \footnotesize\textit{A Level I Specialty card, belonging to Deemer the Warlock. It uses \svg{empower}, but does NOT belong to any School of Magic, despite the imagery.}
}

\columnbreak
\pagetarget{Artifact}{\subheader{Artifact Cards}}

Artifact cards\index{Artifact card} have a top and bottom effect. When you play one, you must \textbf{choose which effect} you are using.
They are divided into 3 Levels: Minor, Major, and Relic, with different border color.
These Levels relate to the overall Power of the card and may be referenced when resolving certain effects or during Scenario setup.
Otherwise, all Artifact cards are normally shuffled together to form the Artifact Deck regardless of their Level.
They are gained through map exploration.\par
Artifacts can be \pagelink{Trading}{traded} in Alliance and Cooperative Scenarios.\par
\begin{tikzpicture}
  \node (major) {\includegraphics[width=0.6\linewidth]{\cards/artifact.png}};
  \begin{relcoords}{major}
    \draw (0.2, 0.8) node {\encircle{1}};
    \draw (0.7, 0.4) node {\encircle{2}};
    \draw (0.9, 0.9) node {\encircle{3}};
    \draw (0.25, 0.2) node {\encircle{4}};
  \end{relcoords}
  \node [below=0pt of major.south] {\imagecaption{Major Artifact}};
  \node (relic) [right=0.05\linewidth of major.south east, anchor=south west] {
\includegraphics[width=0.3\linewidth]{\cards/relic_artifact.png}};
  \begin{relcoords}{relic}
    \draw (0.5, 0.7) node {\encircle{3}};
    \node [below=0pt of relic.south] {\imagecaption{Relic Artifact}};
  \end{relcoords}
  \node (minor) [above=4em of relic.north] {\includegraphics[width=0.3\linewidth]{\cards/minor_artifact.png}};
  \begin{relcoords}{minor}
    \draw (0.5, 0.7) node {\encircle{3}};
    \node [below=0pt of minor.south] {\imagecaption{Minor Artifact}};
  \end{relcoords}
\end{tikzpicture}

\begin{multicols}{2}
  \footnotesize
  \begin{itemize}
    \item[\textbf{1.}] Name
    \item[\textbf{2.}] Effect
    \columnbreak
    \item[\textbf{3.}] Level
    \item[\textbf{4.}] Flavor text
  \end{itemize}
\end{multicols}

\begin{center}
  \vspace*{\fill}
  {\transparent{0.2}\includegraphics[width=0.9\linewidth]{\art/diplomats_ring.png}}
  \vspace*{\fill}
\end{center}

\pagebreak

\begin{expansion}{stronghold}
\pagetarget{Spell Scroll Cards}{\subheader{Spell Scroll Cards}}

During setup shuffle 10 of the Spell Scroll cards into the Artifact Deck.
Set the remaining 10 cards aside to form a Spell Scroll Deck.
Some locations or Scenario effects may call for them later.
\vspace*{1em}

{
    \medskip
    \centering
    \includegraphics[width=0.7\linewidth]{\cards/spell_scroll.png}\\
    \medskip
    \begin{center}
         \imagecaption{Spell Scroll card}
    \end{center}
}

\vspace*{1em}
When you gain a Spell Scroll, place it next to your Hero Board, take the 2 top cards from the Spell Deck and place them face-down on the Spell Scroll.
You can look at these Spells at any time and they don't count toward your Hand \svg{hand} Limit.
You can only have 2 Spell Scrolls at a time.
If you gain a third one, draw the new Spells normally and then decide whether to discard the newly gained Spell Scroll or replace one you already have (including all Spells on it).

\medskip

You can play Spells from a Spell Scroll like any other Spell with the following exceptions:
\begin{itemize}
    \item They don't count towards your Spell Limit.
    You can also play a Spell from your hand during this Combat Round.
    \item You may use the Spell's \textbf{lowest Power} \svg{empower} \textbf{effect}, without needing to pay its cost (e.g. in case of \wikilink{spells/implosion}{Implosion} the cost is 1).
      You can't Empower the Spell with \svg{empower} from any source.
\end{itemize}
\end{expansion}
\columnbreak
\begin{expansion}{stronghold}
    \begin{itemize}
    \item They can't be used to increase the \svg{empower} of another Spell.
    \item You can \pagelink{selling Spell Scrolls}{sell Spell Scrolls} at a Trading Post for 2 \svg{gold} per each Spell on it.
      This works in addition to other effect of visiting the Trading Post.
    \item After using or selling those Spells, the Spell card is \textbf{removed}, not discarded.
\end{itemize}

If the last Spell on a Spell Scroll is used or sold, place the Spell Scroll back in its Deck.

\medskip
Spell Scrolls are not part of your Deck, therefore they are not carried over between Scenarios in a Campaign.
\end{expansion}

\vfill
\begin{expansion}{rampart,cove}
  \pagetarget{War Machines}{\subheader{War Machine Cards}}

  War Machines\index{War Machines} are permanent cards that can be bought at either a \pagelink{Trading Post}{Trading Post} or a \pagelink{War Machine Factory}{War Machine Factory}.
  If you buy one at the Trading Post, \textbf{you cannot use} any of the other normal functions of that Field during that Visit.
  War machines are also more expensive at the Trading Post.

  {
    \bigskip
    \centering
    \begin{tikzpicture}
      \node (card) [inner sep=0] {\includegraphics[width=0.8\linewidth]{\cards/war_machine.png}};
      \begin{relcoords}{card}
        \draw (0.6, 0.92) node {\encircle{1}};
        \draw (0.5, 0.15) node {\encircle{2}};
        \draw (0.3, 0.3) node {\encircle{3}};
        \draw (0.7, 0.3) node {\encircle{4}};
      \end{relcoords}
    \end{tikzpicture}

    \footnotesize
    \imagecaption{War Machine card}
    \begin{itemize}[itemsep=0pt, parsep=5pt, topsep=0pt, partopsep=0pt]
      \item[\textbf{1.}] Name
      \item[\textbf{2.}] Effect
      \item[\textbf{3.}] Cost in a War Machine Factory
      \item[\textbf{4.}] Cost in a Trading Post
    \end{itemize}
  }
\end{expansion}

\end{multicols}

\begin{expansion}*{navalbattles}
  \pagetarget{Empowered Ability Token}{\subheader[notoc]{Empowered Ability Tokens}}
\begin{multicols*}{2}
  \setlength\intextsep{0pt}
  \setlength\columnsep{1em}
  \begin{wrapfigure}{r}{0.2\linewidth}
    
\includegraphics[width=\linewidth]{\images/crown-token.png}
  \end{wrapfigure}
  Empowered Ability cards improve your Hero's Ability cards.
  You can gain them by spending Empowered Ability Tokens.
  You can use one of these Tokens at the moment you gain a new Ability card, regardless of how you obtain it.
  If you use such a Token, discard it, remove the just-gained normal Ability card from the game, and gain the corresponding Empowered Ability card in its stead.
  There's only 1 copy of each Empowered Ability card, so if one player takes a particular Empowered Ability card, they prevent other players from obtaining it.
  If an Empowered Ability card isn't available, you can't spend the Empowered Ability Token and must keep it for later turns.

  You can gain Empowered Ability Tokens only in certain Scenarios (e.~g.~``Island of Fire'' from the Naval Battles Mission Book).\\

  \vspace*{5em}
  \begin{center}
    \transparent{0.2}
\includegraphics[width=\linewidth]{\art/scholar.png}
  \end{center}

  \subsection*{Gaining Empowered Ability\\Example}
  \textit{Bob playing as Crag Hack reaches Level II.
    He performs the Search (2) action and draws two cards from the Ability Deck.
    The cards are Offense and Wisdom.
    He chooses Offense and places the other card on the discard pile.}

  \begin{center}
    \includegraphics[width=0.8\linewidth]{\examples/offense_over_wisdom.png}
  \end{center}

  \textit{Bob happens to have an Empowered Ability Token.
    He may either use it, or save the token for later.
    He decides to use it, so he removes the ``Offense'' Ability card and discards the Empowered Ability Token to take the Empowered ``Offense'' Ability card instead.}

  \begin{center}
    \begin{tikzpicture}
      \draw (-3, 0) node {\includegraphics[width=0.45\linewidth]{\cards/ability-offense.png}};
      \draw (-1.5, -1.9) node[circular drop shadow={shadow scale=0.85, shadow xshift=0.5ex, shadow yshift=-0.5ex, opacity=0.7,
        fill=black, path fading={circle with fuzzy edge 15 percent},
        every shadow}]{
\includegraphics[width=0.2\linewidth]{\images/crown-token.png}};
      \draw (1.5, 0) node {\includegraphics[width=0.45\linewidth]{\cards/ability-offense-empowered.png}};

      \draw (-2.8, -3.5) node {\svg[35]{red_cross}};
      \draw (1.6, -3.5) node {\svg[35]{green_tick}};
    \end{tikzpicture}
  \end{center}
  \vspace*{\fill}
\end{multicols*}
\end{expansion}
